\documentclass{article}
\usepackage{amsmath}
\title{Proyecto de titulación}
\author{Rodrigo Antonio Aldrette Salas}
\date{August 2026}

\begin{document}
\maketitle

\section{Introduction}

Motivación
\begin{itemize}
    \item En Gillingham miden mercado secundario, pero mandan toda la variación que sale del mercado de reparaciones
    a accidentes fijos etc.
    \item Manski dice que precios del mercado secundario dependen de los costos de reparación
    \item Las decisiones de reparación son comúnmente no observadas -> Hu /& Xin
    \item La coyuntura climática y discusiones de obsolescencia programada/planificada,
    right to repair, entre otros, hacen importante estudiar los determinantes de la vida
    util de los vehículos
    \item Es de vital importancia ver el efecto que tiene el mercado de reparación en
    el bolsillo de las personas y el mercado 
\end{itemize}

\section{Modelo}

El modelo tendrá varias consideraciones sobre Gillingham
\begin{itemize}
    \item Una vs dos etapas (matiza el efecto de reparar y de info asimétrica)
    \item Parametrizado vs no parametrizada la regla de transición de $s$
    \item Beliefs vs un precio por nivel de s (mejora estimation time por bajar dimensión
    de P, pero tener que calcular otro punto fijo?)
    \item Varios tipos
    \item Varios coches
    \item Varios niveles de precio de reparación
    
\end{itemize}



\end{document}